\chapter{公式}
    \section{\texorpdfstring{$\sum_{k=1}^n k^2 = n(n+1)(2n+1)/6$}{2乗和の公式}}
        \label{2乗和の公式}
        \begin{proof}
            \quad\par
            $S_n := \sum_{k=1}^n k^2$とする。
            $S_n$に何らかの項を付け加え、$T_n$とし、$T_{n+1} - T_n$が$n$の高々1次式となるようにできれば、$T_n$を求めることができ、$S_n$も求まる。
            $S_{n+1} - S_n = (n+1)^2$であり、$n^2$の項を消去するには$S_n$に$n^3$に比例する項を付け加えればよい($n$の2次式では不足である。
            なぜなら$T_{n+1} - T_n$の段階で$n^2$に比例する項が相殺して消えるから$S_{n+1} - S_n$由来の$n^2$の項を決してキャンセルできないからである)。
            その比例係数を$\alpha$とし、$T_n := S_n + \alpha n^3$とすると次式が成り立つ。
            \[ T_{n+1} - T_n = (n+1)^2 - \alpha\bigl((n+1)^3-n^3\bigr) = (n+1)^2 + \alpha(3n^3+3n+1) \]
            そこで$\alpha=-1/3$とすると次式を得る。
            \[ T_{n+1} - T_n = 2/3 + n \]
            よって$n\geq 2$のとき
            \[ T_n = T_1 + \sum_{k=1}^{n-1} (T_{k+1} - T_k) = 2/3 + (n-1)(n/2 + 2/3) \]
            上式は$n=1$のとき$2/3$となるが、これは$T_1$と等しいので、上式は$n=1$のときも含めて成り立つ。
            以上より
            \[ S_n = T_n + n^3/3 = \frac{1}{6}n(n+1)(2n+1) \]
        \end{proof}
    \section{\texorpdfstring{$\sum_{k=0}^n (2k+1)^2 = (n+1)(2n+1)(2n+3)/3$}{奇数の2乗和の公式}}
        \begin{proof}
            \begin{align*}
                \sum_{k=0}^n (2k+1)^2 &= \sum_{k=0}^n (4k^2+4k+1) = 4\sum_{k=1}^n k^2 + 4\sum_{k=1}^n k + n+1 \\
                &= \frac{2}{3}n(n+1)(2n+1) + 2n(n+1) + (n+1) \quad (\because \ref{2乗和の公式}) \\
                &= \frac{1}{3}(n+1)\left[2n(2n+1) + 6n + 3\right] = \frac{1}{3}(n+1)(2n+1)(2n+3)
            \end{align*}
        \end{proof}
    \section{\texorpdfstring{$1\cdot3\cdot5\dots(2n-1)2n = (2n)!/n!$}{1・3・5・(2n-1)・2n = (2n!)/n!}}
        \begin{proof}
            \begin{align*}
                1 \cdot 3 \cdot 5 \dotsm (2n-1)2n = \frac{1 \cdot 2 \cdot 3 \cdot 4 \cdot 5 \cdot 6 \cdot 7 \dotsm (2n-2)(2n-1)2n}{\underbrace{2 \cdot 4 \cdot 6 \dotsm(2n-2)2n}_{n\text{個}}}2^n = \frac{(2n)!}{1 \cdot 2 \cdot 3 \dotsm (n-1)n} = \frac{(2n)!}{n!}
            \end{align*}
        \end{proof}
    \section{\texorpdfstring{$a^n-b^n = (a-b)\sum_{k=0}^{n-1} a^{n-1-k}b^k$}{a\string^n-b\string^nの分解}}
        \begin{proof}
            \begin{align*}
                (a-b)\sum_{k=0}^{n-1} a^{n-1-k}b^k &= \sum_{k=0}^{n-1} a^{n-k}b^k - \sum_{k=0}^{n-1} a^{n-1-k}b^{k+1} \\
                &= a^n + \sum_{k=1}^{n-1} a^{n-k}b^k - \sum_{k=1}^{n} a^{n-k}b^{k} = a^n - b^n
            \end{align*}
        \end{proof}
    \section{M\"obiusの反転公式}
        \begin{shadebox}
            $f,g$を自然数から複素数への写像とする。
            \begin{enumerate}
                \item 任意の自然数$n$に対して$g(n) = \sum_{d \mid n}f(d)$が成り立つならば$f(n) = \sum_{d \mid n} \mu(d)g(n/d) = \sum_{d \mid n} \mu(n/d)g(n)$が成り立つ。
                \item 任意の自然数$n$に対して$g(n) = \prod_{d \mid n}f(d)$が成り立つならば$f(n) = \prod_{d \mid n} g(n/d)^{\mu(d)} = \prod_{d \mid n} g(d)^{\mu(n/d)}$が成り立つ。
            \end{enumerate}
        \end{shadebox}
        \begin{proof}
            \quad\par
            (1について)
            \par
            結論の式の左から2つ目の等号は約数の意味から明らかである。
            1つ目の等号成立を示す。
            \begin{align*}
                \sum_{d \mid n} \mu(d)g(n/d) &= \sum_{d \mid n} \mu(d)\sum_{e \mid n/d} f(e) = \sum_{d \mid n} \sum_{e \mid n/d} \mu(d)f(e) = \sum_{\genfrac{}{}{0pt}{}{d,e}{de \mid n}} \mu(d)f(e) \\
                &= \sum_{e \mid n} \sum_{d \mid n/e} \mu(d)f(e) = \sum_{e \mid n} f(e) \sum_{d \mid n/e} \mu(d) = \sum_{e \mid n} f(e) \indII{e=n} = f(n)
            \end{align*}
            \newline
            \newline
            (2について)
            \par
            \begin{align*}
                \prod_{d \mid n} g(n/d)^{\mu(d)} &= \prod_{d \mid n} \left(\prod_{e \mid n/d} f(e)\right)^{\mu(d)} = \prod_{d \mid n} \prod_{e \mid n/d} f(e)^{\mu(d)} = \prod_{\genfrac{}{}{0pt}{}{d,e}{de \mid n}} f(e)^{\mu(d)} \\
                &= \prod_{e \mid n} \prod_{d \mid n/e} f(e)^{\mu(d)} = \prod_{e \mid n} f(e)^{\sum_{d \mid n/e} \mu(d)} = \prod_{e \mid n} f(e)^{\indII{e=n}} = f(n)
            \end{align*}
        \end{proof}